\section{Apa itu Kriptografi?}
\textbf{Kriptografi} berasal dari bahasa Yunani: \textit{cryptos} (secret/hidden) dan \textit{graphein} (writing), yang secara harfiah berarti "tulisan rahasia".

Menurut Schneier (1996), kriptografi adalah ilmu dan seni untuk menjaga keamanan pesan. Secara lebih teknis, kriptografi adalah ilmu yang mempelajari teknik-teknik matematika yang berhubungan dengan aspek keamanan informasi seperti kerahasiaan, integritas data, serta otentikasi (Menezes, 1996).

\subsection{Layanan Kriptografi}
Untuk mengatasi masalah keamanan di atas, kriptografi menyediakan empat layanan utama:
\begin{enumerate}
    \item \textbf{Kerahasiaan pesan} (\textit{Confidentiality/privacy/secrecy})
    \item \textbf{Keaslian pesan} (\textit{Data integrity})
    \item \textbf{Keaslian pengirim dan penerima} (\textit{Authentication})
    \item \textbf{Anti penyangkalan} (\textit{Non-repudiation})
\end{enumerate}

